% Part: normal-modal-logic
% Chapter: tableaux
% Section: rules-for-K

\documentclass[../../../include/open-logic-section]{subfiles}

\begin{document}

\olfileid{nml}{tab}{rul}

\olsection{\Ax{K} 的规则}

通常命题联结词的规则与通常的命题带号分析表规则相同，只是添加了前缀。
在每种情况下，应用于带号公式
$\sFmla{S}{!A}[\sigma]$ 的规则所产生的新公式也带有前缀~$\sigma$。
这在直观上应该很清楚：例如，如果 $!A \land !B$ 在（名为）~$\sigma$ 的世界上为真，那么 $!A$ 和 $!B$ 在~$\sigma$ 处为真（而不在其他世界上）。我们将这些命题规则列于 \olref{tab:prop-rules} 中。

\begin{table}
  \[\def\arraystretch{3}\begin{array}{|c|c|}
    \hline
    \AxiomC{\sFmla{\True}{\lnot !A}[\sigma]}
    \RightLabel{\TRule{\True}{\lnot}}
    \UnaryInfC{\sFmla{\False}{!A}[\sigma]}
    \DisplayProof
    &
    \AxiomC{\sFmla{\False}{\lnot !A}[\sigma]}
    \RightLabel{\TRule{\False}{\lnot}}
    \UnaryInfC{\sFmla{\True}{!A}[\sigma]}
    \DisplayProof
    \\[1ex]
    \hline
    \AxiomC{\sFmla{\True}{!A \land !B}[\sigma]}
    \RightLabel{\TRule{\True}{\land}}
    \UnaryInfC{\sFmla{\True}{!A}[\sigma]}
    \noLine
    \UnaryInfC{\sFmla{\True}{!B}[\sigma]}
    \DisplayProof
    &
    \AxiomC{\sFmla{\False}{!A \land !B}[\sigma]}
    \RightLabel{\TRule{\False}{\land}}
    \UnaryInfC{$\sFmla{\False}{!A}[\sigma] \quad \mid \quad
      \sFmla{\False}{!B}[\sigma]$}
    \DisplayProof
    \\[2ex]
    \hline
    \AxiomC{\sFmla{\True}{!A \lor !B}[\sigma]}
    \RightLabel{\TRule{\True}{\lor}}
    \UnaryInfC{$\sFmla{\True}{!A}[\sigma] \quad \mid \quad
      \sFmla{\True}{!B}[\sigma]$}
    \DisplayProof
    &
    \AxiomC{\sFmla{\False}{!A \lor !B}[\sigma]}
    \RightLabel{\TRule{\False}{\lor}}
    \UnaryInfC{\sFmla{\False}{!A}[\sigma]}
    \noLine
    \UnaryInfC{\sFmla{\False}{!B}[\sigma]}
    \DisplayProof
    \\[2ex]
    \hline
    \AxiomC{\sFmla{\True}{!A \lif !B}[\sigma]}
    \RightLabel{\TRule{\True}{\lif}}
    \UnaryInfC{$\sFmla{\False}{!A}[\sigma] \quad \mid
      \quad \sFmla{\True}{!B}[\sigma]$}
    \DisplayProof
    &
    \AxiomC{\sFmla{\False}{!A \lif !B}[\sigma]}
    \RightLabel{\TRule{\False}{\lif}}
    \UnaryInfC{\sFmla{\True}{!A}[\sigma]}
    \noLine
    \UnaryInfC{\sFmla{\False}{!B}[\sigma]}
    \DisplayProof
    \\[2ex]
    \hline
  \end{array}\]
  \caption{命题联结词的前缀分析表规则}
  \ollabel{tab:prop-rules}
\end{table}

闭合条件与普通分析表相同，
只是我们要求不仅公式必须匹配，前缀也必须匹配。
因此，若对某个前缀 $\sigma$ 和公式~$!A$，一个分支同时包含
\[
\sFmla{\True}{!A}[\sigma] \quad\text{和}\quad \sFmla{\False}{!A}[\sigma]
\]，
则该分支闭合。

设置假设的规则也与普通分析表相同，
只不过对于假设，我们总是使用前缀~$1$。
（使用哪个前缀并不重要，只要所有假设使用相同的前缀即可。）
因此，例如，我们说
\[
!B_1, \dots, !B_n \Proves !A
\]
当且仅当存在一个以
\[
\sFmla{\True}{!B_1}[1], \dots, \sFmla{\True}{!B_n}[1],
\sFmla{\False}{!A}[1].
\]
为假设的闭分析表。

对于模态算子\iftag{prvBox}{\iftag{prvDiamond}{~$\Box$ 和}{~$\Box$}}{}\iftag{prvDiamond}{~$\Diamond$}{}，
当规则应用于带有前缀~$\sigma$ 的公式时，
其结论的前缀为 $\sigma.n$。
然而，允许的 $n$ 取决于其符号是~$\True$ 还是~$\False$。

\iftag{prvBox}{规则 $\TRule{\True}{\Box}$ 对包含 $\sFmla{\True}{\Box !A}[\sigma]$ 的分支，通过添加 $\sFmla{\True}{!A}[\sigma.n]$ 进行扩展。\iftag{prvDiamond}{ 类似地，}{}}{}\iftag{prvDiamond}{规则 $\TRule{\False}{\Diamond}$ 对包含 $\sFmla{\False}{\Diamond !A}[\sigma]$ 的分支，通过添加 $\sFmla{\False}{!A}[\sigma.n]$ 进行扩展。}{}
\iftag{notprvBox,notprvDiamond}{它}{它们}只能针对某个在应用该规则的分支上\emph{已经}出现的前缀~$\sigma.n$ 应用。我们称这样的前缀（在该分支上）为“已使用的”。

\iftag{prvBox}{规则 $\TRule{\False}{\Box}$ 对包含 $\sFmla{\False}{\Box !A}[\sigma]$ 的分支，通过添加 $\sFmla{\False}{!A}[\sigma.n]$ 进行扩展。\iftag{prvDiamond}{ 类似地，}{}}{}\iftag{prvDiamond}{规则 $\TRule{\True}{\Diamond}$ 对包含 $\sFmla{\True}{\Diamond !A}[\sigma]$ 的分支，通过添加 $\sFmla{\True}{!A}[\sigma.n]$ 进行扩展。}{}
然而，\iftag{notprvBox,notprvDiamond}{这条规则}{这些规则}只能针对某个在应用该规则的分支上\emph{尚未}出现的前缀~$\sigma.n$ 应用。我们称这样的前缀对该分支是“新的”。

这些规则列于\olref{tab:rules-K}中。

\begin{table}
  \begin{center}
    \def\arraystretch{3}\def\fCenter{}
    \begin{tabular}{|c|c|}
    \hline
    \iftag{prvBox}{
      \AxiomC{\sFmla{\True}{\Box !A}[\sigma]}
      \RightLabel{\TRule{\True}{\Box}}
      \UnaryInfC{\sFmla{\True}{!A}[\sigma.n]}
      \DisplayProof
      &
      \AxiomC{\sFmla{\False}{\Box !A}[\sigma]}
      \RightLabel{\TRule{\False}{\Box}}
      \UnaryInfC{\sFmla{\False}{!A}[\sigma.n]}
      \DisplayProof\\
      $\sigma.n$ is used & $\sigma.n$ is new
      \\[1ex]
      \hline}{}
    \iftag{prvDiamond}{
      \AxiomC{\sFmla{\True}{\Diamond !A}[\sigma]}
      \RightLabel{\TRule{\True}{\Diamond}}
      \UnaryInfC{\sFmla{\True}{!A}[\sigma.n]}
      \DisplayProof
      &
      \AxiomC{\sFmla{\False}{\Diamond !A}[\sigma]}
      \RightLabel{\TRule{\False}{\Diamond}}
      \UnaryInfC{\sFmla{\False}{!A}[\sigma.n]}
      \DisplayProof\\
      $\sigma.n$ is new & $\sigma.n$ is used
      \\[1ex]
      \hline}{}
    \end{tabular}
  \end{center}
  \caption{\Ax{K} 的模态规则}
  \ollabel{tab:rules-K}
\end{table}

要求\iftag{prvBox}{$\TRule{\True}{\Box}$}{$\TRule{\False}{\Diamond}$}的前缀必须是已使用的前缀，这一限制是必要的，否则我们会把下面的分析表也算作闭分析表：
\iftag{notprvBox,notprvDiamond}{%
  \iftag{prvBox}{%
  \begin{oltableau}
    [\pFmla{\True}{\Box \formula{A}}{1}, just = \TAss
      [\pFmla{\False}{\lnot\Box\lnot \formula{A}}{1}, just = \TAss
        [\pFmla{\True}{\formula{A}}{1.1}, just = {\TRule{\True}{\Box}[1]}
          [\pFmla{\True}{\Box\lnot\formula{A}}{1},
            just ={\TRule{\False}{\lnot}[2]}
            [\pFmla{\True}{\lnot\formula{A}}{1.1},
              just = {\TRule{\True}{\Box}[4]}
              [\pFmla{\False}{\formula{A}}{1.1},
                  just ={\TRule{\True}{\lnot}[5]}, close]
            ]
          ]
        ]
      ]
    ]
  \end{oltableau}
  }{
  \begin{oltableau}
    [\pFmla{\True}{\lnot\Diamond\lnot \formula{A}}{1}, just = \TAss
      [\pFmla{\False}{\Diamond\formula{A}}{1}, just = \TAss
        [\pFmla{\False}{\formula{A}}{1.1}, just = {\TRule{\False}{\Diamond}[2]}
          [\pFmla{\False}{\Diamond\lnot\formula{A}}{1},
            just ={\TRule{\True}{\lnot}[1]}
            [\pFmla{\False}{\lnot\formula{A}}{1.1},
              just = {\TRule{\False}{\Diamond}[4]}
              [\pFmla{\True}{\formula{A}}{1.1},
                  just ={\TRule{\True}{\lnot}[5]}, close]
            ]
          ]
        ]
      ]
    ]
  \end{oltableau}
}}{
  \begin{oltableau}
    [\pFmla{\True}{\Box \formula{A}}{1}, just = \TAss
      [\pFmla{\False}{\Diamond \formula{A}}{1}, just = \TAss
        [\pFmla{\True}{\formula{A}}{1.1}, just = {\TRule{\True}{\Box}[1]}
          [\pFmla{\False}{\formula{A}}{1.1},
            just = {\TRule{\False}{\Diamond}[2]}, close]
        ]
      ]
    ]
\end{oltableau}}

但是$\Box \formula{A} \Entails/ \Diamond \formula{A}$，因此我们的证明系统将不可靠。类似地，$\Diamond \formula{A} \Entails/
\Box \formula{A}$，但若去掉要求
\iftag{prvBox}{\TRule{\False}{\Box}}{\TRule{\True}{\Diamond}}的前缀必须是
新前缀这一限制，这就会是一个闭分析表： \iftag{notprvBox,notprvDiamond}{%
  \iftag{prvBox}{%
  \begin{oltableau}
    [\pFmla{\True}{\lnot\Box\lnot \formula{A}}{1}, just = \TAss
      [\pFmla{\False}{\Box \formula{A}}{1}, just = \TAss
        [\pFmla{\False}{\formula{A}}{1.1}, just = {\TRule{\True}{\Box}[2]}
          [\pFmla{\False}{\Box\lnot\formula{A}}{1},
            just ={\TRule{\True}{\lnot}[1]}
            [\pFmla{\False}{\lnot\formula{A}}{1.1},
              just = {\TRule{\False}{\Box}[4]}
              [\pFmla{\True}{\formula{A}}{1.1},
                  just ={\TRule{\False}{\lnot}[5]}, close]
            ]
          ]
        ]
      ]
    ]
  \end{oltableau}
  }{
  \begin{oltableau}
    [\pFmla{\True}{\Diamond \formula{A}}{1}, just = \TAss
      [\pFmla{\False}{\lnot\Diamond\lnot\formula{A}}{1}, just = \TAss
        [\pFmla{\True}{\formula{A}}{1.1}, just = {\TRule{\True}{\Diamond}[1]}
          [\pFmla{\True}{\Diamond\lnot\formula{A}}{1},
            just ={\TRule{\False}{\lnot}[2]}
            [\pFmla{\True}{\lnot\formula{A}}{1.1},
              just = {\TRule{\True}{\Diamond}[4]}
              [\pFmla{\False}{\formula{A}}{1.1},
                  just ={\TRule{\True}{\lnot}[5]}, close]
            ]
          ]
        ]
      ]
    ]
  \end{oltableau}
}}{
  \begin{oltableau}
    [\pFmla{\True}{\Diamond \formula{A}}{1}, just = \TAss,name=one
      [\pFmla{\False}{\Box \formula{A}}{1}, just = \TAss, name=two
        [\pFmla{\True}{\formula{A}}{1.1}, just = {\TRule{\True}{\Diamond}[1]}
          [\pFmla{\False}{\formula{A}}{1.1},
            just = {\TRule{\False}{\Box}[2]}, close]
        ]
      ]
    ]
  \end{oltableau}}

\end{document}